\begin{table}[H]
  \centering
  \small
  \caption{Estrategias de fragmentación consideradas (ADR-0001).}
  \label{tab:alternativas_chunking}
  \begin{tabularx}{\textwidth}{@{}>{\raggedright\arraybackslash}p{3.2cm} X@{}}
    \toprule
    \textbf{Estrategia} & \textbf{Descripción del método} \\
    \midrule
    Tamaño fijo con solape &
    Corta el texto cada $n$ tokens con ventana de solape. Simple, pero rompe frases y puede mezclar dos asignaturas en un mismo fragmento. \\
    \addlinespace
    Estructural por unidad semántica &
    Parte de la unidad docente (guía o bloque de salidas) y solo subdivide si excede el máximo, cortando por fronteras naturales. Respeta los límites entre asignaturas por construcción. \\
    \addlinespace
    Semántica por incrustaciones &
    Detecta cambios de tema comparando \emph{embeddings} de frases consecutivas.\\
    \addlinespace
    Un fragmento por asignatura &
    Sin subdivisión. El 50\,\% de las guías supera los 2\,677 caracteres (máx.\ 23\,940), por lo que la incrustación truncaría contenido en silencio.\\
    \bottomrule
  \end{tabularx}
\end{table}

